\section{Detection of global powers at degree-one primes}
\label{sd-section}

We now prove the power-detection statement used in
assertion~\ref{sc-few}.  The argument applies to every positive
integer exponent.  Its classical input is the following
qualitative form of the Chebotarev theorem: in a finite
Galois extension of \(\mathbb Q\), every conjugacy class
occurs as the Frobenius class of infinitely many unramified
rational primes.\cite[Theorem 8.31]{milne-ant}
No effective prime bound is needed.

A prime ideal \(P\) of a number field \(F\) has
\emph{absolute residue degree one} if its residue field
is \(\mathbb F_r\), where \(r\) is the rational prime
below \(P\).
For \(\beta\in F^\times\), reduction modulo \(P\) is
defined as a nonzero residue when \(v_P(\beta)=0\).
Equivalently, there are \(a,b\in\mathcal O_F\) with
\(b\notin P\), \(\beta b=a\), and \(a\notin P\).
The residue is then \(\overline a/\overline b\).
This is precisely the residue relation used by the source:
the definition of the local ring
\(\mathcal O_{F,P}\) gives the converse as well.

\begin{lemma}\label{sd-degree-one}
Let \(K/F\) be a finite Galois extension of number fields,
let \(\sigma\in\operatorname{Gal}(K/F)\), and let
\(\mathcal B\) be any finite set of primes of \(F\).
There exists an unramified prime \(P\notin\mathcal B\)
of absolute residue degree one, and a prime of \(K\)
above \(P\), whose arithmetic Frobenius is \(\sigma\).
\end{lemma}
\begin{proof}
Choose a finite Galois extension \(N/\mathbb Q\)
containing \(K\).  The restriction map
\(\operatorname{Gal}(N/F)\to\operatorname{Gal}(K/F)\)
is surjective: an \(F\)-embedding of the subfield \(K\)
extends to an \(F\)-automorphism of the normal
extension \(N/F\).
Choose a lift \(\widetilde\sigma\).

Chebotarev provides infinitely many rational primes
whose Frobenius class in \(\operatorname{Gal}(N/\mathbb Q)\)
is the conjugacy class of \(\widetilde\sigma\).
Omit the finitely many rational primes lying below
\(\mathcal B\).
For any remaining prime \(r\), choose a prime \(Q\) of
\(N\) at which the Frobenius element is exactly
\(\widetilde\sigma\), rather than a different conjugate.
This is possible because changing \(Q\) by an
automorphism conjugates its Frobenius.
Put \(P=Q\cap\mathcal O_F\).

For \(a\in\mathcal O_F\), the element
\(\widetilde\sigma\) fixes \(a\), whereas its action
on the residue field at \(Q\) is \(x\mapsto x^r\).
Every element of \(\mathcal O_F/P\) is therefore
fixed by the \(r\)-power map.  The polynomial
\(X^r-X\) has at most \(r\) roots, and the residue
field contains the \(r\) elements of \(\mathbb F_r\).
Consequently \(\mathcal O_F/P=\mathbb F_r\).
All these primes are unramified in \(N/\mathbb Q\).
Since \(NP=r\), the restriction of
\(\widetilde\sigma\) to \(K\) is also the arithmetic
Frobenius of \(K/F\) at the chosen prime above \(P\).
It is \(\sigma\), as required.
\end{proof}

\begin{lemma}\label{sd-roots-reduce}
Let \(F\) contain all \(q\)-th roots of unity, and let
\(P\nmid q\) be a finite prime.
Reduction modulo \(P\) is injective on \(\mu_q\).
\end{lemma}
\begin{proof}
Every root of unity and its inverse are algebraic integers,
so these reductions are defined and nonzero.
If \(\zeta\in\mu_q\) reduces to \(1\) and \(\zeta\ne1\),
then
\[
             1+\zeta+\cdots+\zeta^{q-1}=0.
\]
Reducing this equality modulo \(P\) gives \(q=0\),
contrary to \(P\nmid q\).
Apply this to the quotient of two roots of unity
with equal reduction.
\end{proof}

\begin{theorem}\label{sd-power}
Let \(F\) be a number field, let \(q\ge1\), and suppose
\(\mu_q\subset F\).
Let \(\beta\in F^\times\) and let \(\mathcal B\) be
a finite set of primes.
Suppose that, at every prime \(P\notin\mathcal B\)
of absolute residue degree one at which \(\beta\)
is a unit, its residue is a \(q\)-th power in
\((\mathcal O_F/P)^\times\).
Then \(\beta\) is a \(q\)-th power in \(F^\times\).
\end{theorem}
\begin{proof}
The case \(q=1\) is immediate.
For \(q>1\), choose \(\alpha\) with
\(\alpha^q=\beta\) and put \(K=F(\alpha)\).
The polynomial is separable and all its roots
are \(\zeta\alpha\), \(\zeta\in\mu_q\subset F\).
Thus \(K/F\) is Galois.

If \(\alpha\notin F\), choose
\(\sigma\in\operatorname{Gal}(K/F)\) with
\(\sigma\alpha\ne\alpha\).
Enlarge \(\mathcal B\) by the primes dividing \(q\)
and by the finite support of the principal divisor
of \(\beta\).
Lemma~\ref{sd-degree-one} gives a prime \(P\) and a
prime \(Q\) of \(K\) above it with Frobenius \(\sigma\).
Let \(r=NP\), a rational prime.
The equality \(\alpha^q=\beta\) shows that \(\alpha\)
is a unit at \(Q\).
By the hypothesis there is \(t\in\mathbb F_r^\times\)
with \(t^q=\overline\beta\).

All \(q\)-th roots of unity reduce to elements of
\(\mathbb F_r\), and they remain distinct by
Lemma~\ref{sd-roots-reduce}.
The \(q\) roots of \(X^q-\overline\beta\) in any
extension of the residue field are consequently the
elements \(\overline\zeta t\), all in \(\mathbb F_r\).
In particular \(\overline\alpha\in\mathbb F_r^\times\).
The Frobenius property gives
\[
             \overline{\sigma\alpha}
                  =\overline\alpha^{\,r}
                  =\overline\alpha.
\]
But \(\sigma\alpha/\alpha\in\mu_q\), so its reduction
is \(1\).  The same injectivity lemma forces
\(\sigma\alpha/\alpha=1\), a contradiction.
Hence \(\alpha\in F^\times\), proving the theorem.
\end{proof}

\begin{corollary}\label{sd-source}
Let \(L\subset F\) be number fields, let
\(\mu_q\subset F\), let \(\beta\in F^\times\),
and let \(0\ne b\in\mathcal O_F\).
Assume that for every prime \(P\nmid b\) such that
the natural map
\[
                 \mathcal O_L\longrightarrow\mathcal O_F/P
\]
is surjective, \(\beta\) has a nonzero residue
which is a \(q\)-th power.
Then \(\beta\in(F^\times)^q\).
In particular, this applies to \(q=\ell^\nu\).
\end{corollary}
\begin{proof}
At every absolute degree-one prime \(P\), the residue
field is \(\mathbb F_r\).  The image of \(\mathcal O_L\)
contains the image of \(\mathbb Z\), which already
equals \(\mathbb F_r\); the displayed map is therefore
surjective.  Apply Theorem~\ref{sd-power}, with the
primes dividing \(b\) excluded.
No normality assumption on \(F/L\) is required.
\end{proof}

The roots-of-unity hypothesis is essential to this argument
and to the resulting statement.  A precise counterexample
after deleting it is
\[
                      F=\mathbb Q,\qquad q=8,\qquad\beta=16.
\]
The element \(16\) is not an eighth power in \(\mathbb Q\),
because its \(2\)-adic valuation is \(4\).
Nevertheless it is an eighth power modulo every odd
rational prime \(r\).
Indeed at least one of \(2,-2,-1\) is a square in
\(\mathbb F_r^\times\): the square subgroup has index
two, and their product \(4\) is a square, so they cannot
all belong to the other coset.
If \(t^2=2\) or \(t^2=-2\), then \(t^8=16\).
If \(i^2=-1\), then \((1+i)^2=2i\), and hence
\((1+i)^8=16\).
All these elements are nonzero when \(r\) is odd.
This example does not satisfy \(\mu_8\subset\mathbb Q\),
which is an explicit hypothesis of the source.

The detection step thus has a complete conventional proof,
with Chebotarev as the stated classical input.
It supplies the final separation argument for the finite-set
comparison; the finite-field homology classes and the integral
cycles whose values must be separated are constructed separately.


